% p173_1; source PDF page 173; manually verified editable table.
\begin{tikzpicture}[x=1cm,y=-1cm,inner sep=0pt,outer sep=0pt]
\path[use as bounding box] (0,0) rectangle (16.000,8.920);
\draw[line width=.4pt] (0.000,0.000) rectangle (4.100,0.805);
\node[align=center,text width=3.840cm,font={\linespread{1}\fontsize{10}{13.50}\selectfont\bfseries}] at (2.050,0.403) {罪名关系的类型};
\draw[line width=.4pt] (4.100,0.000) rectangle (10.300,0.805);
\node[align=center,text width=5.940cm,font={\linespread{1}\fontsize{10}{13.50}\selectfont\bfseries}] at (7.200,0.403) {特点};
\draw[line width=.4pt] (10.300,0.000) rectangle (16.000,0.805);
\node[align=center,text width=5.440cm,font={\linespread{1}\fontsize{10}{13.50}\selectfont\bfseries}] at (13.150,0.403) {范例};
\draw[line width=.4pt] (0.000,0.805) rectangle (4.100,3.610);
\node[align=left,text width=3.840cm,font={\linespread{1}\fontsize{10}{13.50}\selectfont}] at (2.050,2.208) {A 与 $-A$ 的对立排斥关系（男生与女生的逻辑关系）};
\draw[line width=.4pt] (4.100,0.805) rectangle (10.300,3.610);
\node[align=left,text width=5.940cm,font={\linespread{1}\fontsize{10}{13.50}\selectfont}] at (7.200,2.208) {一个行为针对一个对象，不可能同时触犯这两个罪名，也即既不可能属于法条竞合，也不可能形成想象竞合};
\draw[line width=.4pt] (10.300,0.805) rectangle (16.000,3.610);
\node[align=left,text width=5.440cm,font={\linespread{1}\fontsize{10}{13.50}\selectfont}] at (13.150,2.208) {甲欺骗乙回头看看，甲趁乙回头，拿走乙的财物。甲对乙，要么定诈骗罪，要么定盗窃罪；应定盗窃罪};
\draw[line width=.4pt] (0.000,3.610) rectangle (4.100,6.415);
\node[align=left,text width=3.840cm,font={\linespread{1}\fontsize{10}{13.50}\selectfont}] at (2.050,5.013) {A 与 A+B 的包容评价关系（本科生与硕士生的逻辑关系）};
\draw[line width=.4pt] (4.100,3.610) rectangle (10.300,6.415);
\node[align=left,text width=5.940cm,font={\linespread{1}\fontsize{10}{13.50}\selectfont}] at (7.200,5.013) {A+B 可以包容评价为 A。触犯 A+B，必然触犯 A。所以这种关系也是法条竞合关系，也被称为位阶关系};
\draw[line width=.4pt] (10.300,3.610) rectangle (16.000,6.415);
\node[align=left,text width=5.440cm,font={\linespread{1}\fontsize{10}{13.50}\selectfont}] at (13.150,5.013) {甲想盗窃普通财物，实际盗窃到枪支。盗窃枪支罪可以包容评价为盗窃罪。甲构成盗窃罪既遂};
\draw[line width=.4pt] (0.000,6.415) rectangle (4.100,8.920);
\node[align=left,text width=3.840cm,font={\linespread{1}\fontsize{10}{13.50}\selectfont}] at (2.050,7.668) {A 与 B 的中立关系（男生与硕士生的逻辑关系）};
\draw[line width=.4pt] (4.100,6.415) rectangle (10.300,8.920);
\node[align=left,text width=5.940cm,font={\linespread{1}\fontsize{10}{13.50}\selectfont}] at (7.200,7.668) {A 与 B 不可能属于法条竞合关系，但有可能形成想象竞合};
\draw[line width=.4pt] (10.300,6.415) rectangle (16.000,8.920);
\node[align=left,text width=5.440cm,font={\linespread{1}\fontsize{10}{13.50}\selectfont}] at (13.150,7.668) {上文“偷救命药案”，甲同时构成盗窃罪和故意杀人罪，想象竞合，择一重罪论处，定故意杀人罪};
\end{tikzpicture}
